\section{Sobolev extension across removable dust}\label{sec:sobolev}

Throughout this section, $B\subset\C$ is compact and totally
disconnected, and $G=\Chat\setminus B$.  For a finite vector of scalar
functions, the gradient norm means the $L^2$ norm of its matrix of
gradients.  Small changes of an old test will always mean small changes
both uniformly on the sphere and in gradient norm on $G$.

\begin{theorem}[Continuous Sobolev extension]\label{thm:dust-sobolev}
If $B\subset\C$ is compact, totally disconnected, and conformally
removable, then every continuous real function $u$ on $\Chat$ whose
restriction to $G=\Chat\setminus B$ belongs to $W^{1,2}(G)$ belongs
to $W^{1,2}(\Chat)$.
\end{theorem}

The next lemma supplies small-energy barriers around the source ends.
Its stability assertion will let us retain those barriers while adding
further tests.

\begin{lemma}[Appending separating tests]\label{lem:append-tests}
There is a radius $D_0$, depending only on $B$, with the following
property.  Suppose $u_1,\ldots,u_m$ satisfy the hypotheses of
Proposition~\ref{prop:models}.  Given $\delta,\eta>0$, they may be
replaced by tests $\widetilde u_i$ satisfying
\[
 \|\widetilde u_i-u_i\|_\infty<\delta,
 \qquad \|\nabla(\widetilde u_i-u_i)\|_{2,G}<\delta,
\]
and two tests $p,q$ may be appended, so that the enlarged finite list
still satisfies the pair-null condition~\eqref{eq:pair-null}.
There are finitely many smooth Jordan curves $\gamma_l\subset G\cap
D_{D_0}$ whose bounded sides cover $B$, and smooth functions
$\Phi_l:\R^2\to[0,1]$, such that
\begin{equation}\label{eq:saved-barriers}
 \begin{gathered}
 \Phi_l(p,q)=1\text{ near }\gamma_l,\qquad
 \Phi_l(p,q)=0\text{ on a fixed collar of }\partial D_{D_0},\\
 \|\nabla\Phi_l(p,q)\|_{2,G\cap D_{D_0}}<\eta.
 \end{gathered}
\end{equation}
The barrier conclusions~\eqref{eq:saved-barriers}, with the same curves
and functions $\Phi_l$, survive sufficiently small subsequent changes
of the coordinates $p,q$.
\end{lemma}

\begin{proof}
We construct a continuous pair $P=(p,q)$ and obtain the barriers by
composing it with translates of a radial cutoff of small Dirichlet
energy.  By conformal invariance, composing the cutoff with a model map
$h$ does not increase its Dirichlet energy, so we will compare $P$ with
$(\Re h,\Im h)$ on a fixed source disk.  Proposition~\ref{prop:height}
supplies a height $v$ whose pullback gives the vertical coordinate.
We make small range compressions of this pullback and the old tests so
that they become constant near every source end corresponding to a
nonpoint model hole, then modify the horizontal coordinate there.
These modifications preserve pair-nullity and make the gradient error
small on the source disk.  The vertical error is uniformly small; the
horizontal changes are confined to neighborhoods where $h$ has small
energy.

Choose $a\ge1$ with $B\subset D_a$.  For every normalized
univalent map $h$ on $G$, the Exterior Area Theorem gives
\begin{equation}\label{eq:sob-exterior}
 |h(z)-z|\le a\sqrt{-\log(1-a^2/|z|^2)}\qquad(|z|>a).
\end{equation}
Indeed, apply Cauchy--Schwarz to the Laurent series and the area-theorem
coefficient bound.  Hence $h(\partial D_{2a})\subset D_{3a}$, and the
omitted set $F=\C\setminus h(G)$ lies on its bounded side.  In particular,
$F\subset D_{3a}$ for every model of every list.  Fix $D_0=64a$.
The image of $G\cap D_{65a}$ lies in $D_{67a}$; these uniform bounds also
bound the energy of $(\Re h,\Im h)$ on that source disk.

Choose a smooth radial cutoff $\Phi$, equal to one on a disk of radius
$r>0$ and zero outside $D_{4a}$, with
$\|\nabla\Phi\|_{2,\C}<\eta/3$.  Such a cutoff is obtained by smoothing
a logarithmic profile and taking $r$ sufficiently small.  All its
translates have the same derivative bounds.  We shall use translates
centered at the eventual traces of $(p,q)$ on $B$.

By Proposition~\ref{prop:height}, take a model $h$ sufficiently close to
the supremum of its real Laurent coefficient that the associated height
$v$ satisfies
\begin{equation}\label{eq:height-errors-for-tests}
 \|v-Y\|_{\infty,D_{67a}}\ \hbox{and}\
 \|\nabla(v-Y)\|_{2,\C}
 \quad\hbox{as small as required below},
 \qquad Y(w)=\Im w.
\end{equation}
Write $U_i$ for the global model continuation of $u_i$.  The function
$v$ is constant on every component of $F$, so $v\circ h$ extends
continuously to $B$ by Lemma~\ref{lem:ends}.

\emph{Flattening the hole values.}
For a continuous, locally Sobolev real function $s$ of finite Dirichlet
energy on a region, the measure
\[
 \mu_s(I)=\int_{\{s\in I\}}|\nabla s|^2
\]
has no atom at any real number: the Sobolev gradient is zero almost
everywhere on each fixed level set.  Thus, given countably many values,
we can surround them by open intervals whose total length and total
$\mu_s$-measure are arbitrarily small.  If $O$ is their union, then
\[
 T(t)=\int_0^t\mathbf1_{\R\setminus O}(\tau)\,\dd\tau
\]
is $1$-Lipschitz, satisfies $\|T-\id\|_\infty\le |O|$, and is constant
near each specified value.  The chain rule bounds the squared gradient
error by $\mu_s(O)$.  This remains valid when the specified values are
dense.

Apply this construction to each $u_i$, at its values at the source ends
corresponding to nonpoint components of $F$, and to $v$ at its values
on those components.  For $v$ use only the energy on $D_{67a}$.
Denote the flattened old tests by $\widetilde u_i=T_i\circ u_i$, and put
$q_0=T\circ v\circ h$.  Their errors can be arbitrarily small, and
the required bounds by $\delta$ hold for the old tests.  In particular,
$q_0$ is uniformly and energetically close to $\Im h$ on $G\cap D_{65a}$.
By Lemma~\ref{lem:level-pairs} and coordinatewise Lipschitz composition,
the traces of $q_0$ form null-area pairs with the flattened old tests.
The old pairs also remain null.  Each nonpoint-hole end $b_j$ now has
a source neighborhood on which all $\widetilde u_i$ and $q_0$ are
constant.

\emph{The horizontal coordinate.}
Let $F_j$ be the nonpoint component corresponding to $b_j$, and choose
$c_j\in\Re F_j$.  Around each $b_j$ choose a ball $B_j\subset D_{2a}$
in the preceding common constancy neighborhood.  Its radius may be
required to tend to zero and its image in $G$ to lie within $2^{-j}a$
of $F_j$.  The balls can also satisfy
\begin{equation}\label{eq:small-change-support}
 \|\mathbf1_S\nabla H\|_2<\tau,
 \qquad S=\bigcup_jB_j,\qquad H=(\Re h,\Im h),
\end{equation}
for any prescribed $\tau>0$: the energy is integrable, and each center
is a single source point.  Choose logarithmic cutoffs $\chi_j$ supported
in $B_j$, equal to one near $b_j$, with
$\sum_j\|\nabla\chi_j\|_2<\tau$.  Put
\begin{equation}\label{eq:priority-interpolation}
 w_j=\chi_j\prod_{i<j}(1-\chi_i),\qquad
 \rho=\prod_j(1-\chi_j),\qquad
 p_0=\rho\Re h+\sum_jw_jc_j\quad\hbox{on }G.
\end{equation}
These expressions are locally finite in $G$, since the radii tend to
zero and all centers belong to $B$.  They are convex interpolations.
On $S$, all horizontal values used have absolute value at most $4a$.
The product rule, applied to the finite formula locally, gives
\begin{equation}\label{eq:priority-gradient}
 |\nabla p_0-\rho\nabla\Re h|
 \le 8a\sum_j|\nabla\chi_j|,\qquad
 \|\nabla(p_0-\Re h)\|_{2,G\cap D_{65a}}
 \le(1+8a)\tau.
\end{equation}
No disjointness of the balls is required.

At $b_j$, the factor $\chi_j=1$ makes~\eqref{eq:priority-interpolation}
terminate after index $j$ on a neighborhood.  It is therefore a finite
continuous expression there.  At a point end $b$, write
$F_b=\{\zeta_b\}$ and $x_b=\Re\zeta_b$.  Early nonzero weights may
change the trace from $x_b$.  To prove continuity, compare $p_0$ with
the finite continuous function
\[
 x_b+\sum_{j\le N}w_j(z)(c_j-x_b).
\]
For sufficiently large $N$ and sufficiently small neighborhoods of $b$,
every later ball meeting that neighborhood has its center near $b$.
Upper semicontinuity of the end correspondence then puts its entire
component $F_j$, and hence $c_j$, arbitrarily close to $\zeta_b$ and
$x_b$.  Also $\Re h(z)\to x_b$ from $G$.  The remaining weight is a
convex combination of these close values, so the difference from the
displayed finite expression is uniformly arbitrarily small.  This proves
continuous extension at $b$, including when early supports overlap.
The same estimate controls the prescribed values at nearby ends.

The new pairs are now checked on all of $B$.  On each ball $B_j$, every
coordinate other than $p_0$ is constant, so a pair involving $p_0$ has
image in a line.  Countably many balls contribute zero planar area.
Outside $S$, all ends are point ends and $p_0$ is their horizontal
trace.  Apply the horizontal assertion of Lemma~\ref{lem:level-pairs}
to $T_i\circ U_i$ and to $T\circ v$, respectively.  Their gradients
vanish on $F$, so this proves nullity of the pairs with every old test
and with $q_0$.  Together with the earlier pair checks, this proves the
full pair-null condition.  Multiply $p_0,q_0$ by a smooth source cutoff
equal to one on $D_{65a}$ and zero outside a larger disk.  The resulting
$p,q$ are bounded continuous tests of finite energy on $G$; no trace or
working-disk calculation changes.

\emph{Small energy and strict margins.}
Let $P=(p,q)$ on the working disk.  Outside $S$ its first coordinate
equals that of $H$, and write $e=\|q-\Im h\|_\infty$ there.  For any
translate $\Phi_l$ of $\Phi$, with derivative bounds
$M_1=\|D\Phi\|_\infty$ and $M_2=\|D^2\Phi\|_\infty$, the chain rule
gives
\begin{align}\label{eq:barrier-comparison}
 \|\nabla\Phi_l(P)-\nabla\Phi_l(H)\|_2
 &\le M_1\|\nabla P-\nabla H\|_2
       +M_2 e\|\nabla H\|_2
       +2M_1\|\mathbf1_S\nabla H\|_2 .
\end{align}
All norms here are on $G\cap D_{D_0}$.  The first and last terms tend
to zero by the choices above; the middle term does so by
\eqref{eq:height-errors-for-tests} and flattening.  This estimate does
not require a small uniform horizontal error on $S$.  Conformal
invariance gives $\|\nabla\Phi_l(H)\|_2<\eta/3$, so the choices can
ensure the strict bound in~\eqref{eq:saved-barriers}, uniformly over
all translates.

Take the total vertical uniform error less than $a$.  On $B$,
$|p|\le3a$ by convexity and $|q|\le4a$, so every trace center has norm
at most $5a$.  On the collar $63a\le |z|\le65a$, the first coordinate
is $\Re h$ and the second differs from $\Im h$ by less than $a$.
Equation~\eqref{eq:sob-exterior} puts its distance from every trace
center above $50a$, well beyond the support radius $4a$.  At each
$b\in B$, the translate centered at $P(b)$ is one wherever
$|P-P(b)|<r$.  Choose a smaller source neighborhood whose closure has
$|P-P(b)|<r/2$.  Compactness and the contour construction in the proof
of Lemma~\ref{lem:ends} give finitely many inner contours in such
neighborhoods with bounded sides covering $B$.

Finally, for a fixed saved pair $A=(p,q)$ and any perturbed pair $A'$,
\begin{equation}\label{eq:saved-stability}
 \|\nabla\Phi_l(A')-\nabla\Phi_l(A)\|_2
 \le M_1\|\nabla A'-\nabla A\|_2
   +M_2\|A'-A\|_\infty\|\nabla A\|_2.
\end{equation}
The positive plateau and zero-region margins preserve the exact values
one and zero under small uniform changes.  The positive energy slack
and~\eqref{eq:saved-stability} preserve the strict norm bound.  This
proves the stability assertion and the lemma.
\end{proof}



\begin{proof}[Proof of Theorem~\ref{thm:dust-sobolev}]
The empty case is immediate.  Fix $0<\epsilon<1$ and start with the
one-element list $u$, for which pair-nullity is vacuous.  Apply
Lemma~\ref{lem:append-tests} successively with $\eta=1/k$, saving at
stage $k$ its pair, curves, and functions $\Phi_{k,l}$.  By the strict
margins and~\eqref{eq:saved-stability}, there is $t_k>0$ such that total
uniform and gradient changes of that pair below $t_k$ preserve all its
saved conclusions with norm less than $1/k$.

At stage $n$, require each earlier scalar coordinate to change by less
than
\[
 d_n=\frac{2^{-n-3}}{\sqrt2}
       \min(1,\epsilon,t_1,\ldots,t_{n-1})
\]
in both norms (omit the $t$'s when $n=1$).  Each stage has only finitely
many requirements, and the lemma permits arbitrarily small changes.
For every fixed $k$, the sum of subsequent pair errors is less than
$t_k/4$.  Thus all saved barriers work at every later finite stage.
If $a_n$ is the first coordinate at stage $n$, then
\begin{equation}\label{eq:first-test-control}
 a_n\longrightarrow u_\epsilon\text{ uniformly},\qquad
 \|u_\epsilon-u\|_\infty<\epsilon,\qquad
 \|\nabla a_n\|_{2,G}\le\|\nabla u\|_{2,G}+1.
\end{equation}
Only finite lists are required to be pair-null.  In particular, no
infinite-list model assertion or limiting pair-null assertion is being
used.

Choose a model $h_n$ of each resulting finite list by
Proposition~\ref{prop:models}.  Lemma~\ref{lem:normalization} gives a
subsequence converging locally uniformly on $G$ to a normalized
univalent map $h$.  We claim that every component of its omitted set is
a point.  Otherwise, by Lemma~\ref{lem:ends}, some source end $b$ has
a nondegenerate target component $K$.  One axial projection of $K$
contains an interval $J$ of positive length.

For each $k$, choose a saved inner contour $\gamma_{k,l}$ enclosing
$b$.  Its limiting image strictly encloses $K$.  Uniform convergence
on the contour and winding-number stability imply that
$h_n(\gamma_{k,l})$ also encloses $K$ for all sufficiently large $n$
in the subsequence.  Use the two continued model coordinates at stage
$n$ in $\Phi_{k,l}$, on the bounded side of
$h_n(\partial D_{D_0})$, and extend by zero outside.  Denote this
function by $V_{n,k}$.  It is continuous and globally Sobolev: the
saved zero collar removes any interface jump, and the model
coordinates are globally Sobolev.  Their gradients vanish on the
model holes.  The chain rule and conformal invariance therefore give
\begin{equation}\label{eq:vanishing-barrier-energy}
 \|\nabla V_{n,k}\|_{2,\C}<1/k.
\end{equation}
It equals one on $h_n(\gamma_{k,l})$.  Moreover all such functions are
zero outside one fixed disk $D_L$, for example with $L=67D_0/64$ by
the uniform bounds in the preceding proof.  The radius does not depend
on $k$, $n$, or the coordinate list.

\begin{figure}[htbp]
\centering
\begin{tikzpicture}[x=1.05cm,y=1.05cm,>=Latex,font=\small]
  \fill[blue!3] (0,0) circle (2.45);
  \draw[blue!55!black,thick] (0,0) circle (2.45);
  \node[anchor=south west] at (1.65,1.78) {$D_L$};
  \draw[thick] (-.40,0) ellipse (1.20 and 1.40);
  \node at (-.69,1.60) {$C$};
  \node at (0,1.98) {$V=1$ on $C$};
  \filldraw[fill=gray!25,draw=gray!70!black]
    (-.89,-.69) .. controls (-1.12,-.18) and (-.61,.03) .. (-.93,.54)
    .. controls (-.55,.82) and (-.08,.49) .. (.07,.70)
    .. controls (.47,.40) and (.02,.09) .. (.34,-.25)
    .. controls (.14,-.79) and (-.44,-.52) .. (-.89,-.69);
  \node at (-.46,-.17) {$K$};
  \draw[gray,densely dashed] (-2.78,-.55) -- (2.75,-.55);
  \draw[gray,densely dashed] (-2.78,.45) -- (2.75,.45);
  \draw[very thick,blue!70!black] (-2.82,-.55) -- (-2.82,.45);
  \draw[blue!70!black] (-2.89,-.55)--(-2.75,-.55);
  \draw[blue!70!black] (-2.89,.45)--(-2.75,.45);
  \node[anchor=east] at (-2.94,-.05) {$J$};
  \draw[gray,densely dashed] (-2.75,0) -- (2.85,0);
  \fill (.80,0) circle (.035);
  \fill (2.45,0) circle (.035);
  \draw[->,very thick,orange!85!black] (.80,0) -- (2.45,0);
  \node[above,orange!75!black] at (1.63,.09) {$1\longrightarrow0$};
  \node[anchor=north west] at (2.48,-.06) {$V=0$};
  \node[align=center] at (0,-2.93)
    {$\displaystyle \int_{\text{slice}}|\nabla V|^2\,\dd s
       \ge \frac{1}{2L}$};
\end{tikzpicture}
\caption{A fixed interval $J$ in a projection of a nondegenerate
component $K$ forces a uniform energy cost. Each corresponding line
meets the inner contour $C$, where the barrier is one, and then reaches
zero within a segment of length at most $2L$. Integrating over $J$
contradicts arbitrarily small barrier energy.
Here $V=V_{n,k}$ and $C=h_n(\gamma_{k,l})$.}
\label{fig:line-payment}
\end{figure}


For definiteness suppose $J$ is the vertical projection interval.
Every horizontal line at a height in $J$ meets the inner image contour, where
$V_{n,k}=1$.  On almost every such line, the Sobolev restriction is
absolutely continuous and agrees with the continuous restriction.
It takes value zero at the left and right ends of the segment with
horizontal coordinates $-L$ and $L$.  Cauchy--Schwarz, followed by
integration over $J$, gives
\begin{equation}\label{eq:line-payment-sobolev}
 \int_\C|\nabla V_{n,k}|^2
 \ge\frac{|J|}{2L}>0.
\end{equation}
The argument only uses the value one on the contour, as illustrated
in Figure~\ref{fig:line-payment}.  Choose $k$ so large that
$k^{-2}<|J|/(2L)$, and then choose $n$ sufficiently late for that
saved contour.  Equations~\eqref{eq:vanishing-barrier-energy}
and~\eqref{eq:line-payment-sobolev} contradict each other.

All limit components are therefore singletons.
Lemma~\ref{lem:ends} now extends $h$ to an orientation-preserving
homeomorphism of the sphere.  Removability of $B$ makes it a M\"obius
map; the normalization $h(z)=z+O(1/z)$ makes it the identity.

Let $A_n$ denote the globally continued first test of the model $h_n$.
Its values on holes are the source endpoint values, and its gradient
vanishes there.  Consequently
\[
 \|A_n\|_\infty\le\|u\|_\infty+1,\qquad
 \|\nabla A_n\|_{2,\Chat}
   =\|\nabla a_n\|_{2,G}\le\|\nabla u\|_{2,G}+1.
\]
Take a weakly convergent subsequence in $W^{1,2}(\Chat)$.  Since
$h_n\to\id$ locally uniformly on $G$, inverse convergence and
\eqref{eq:first-test-control} give $A_n\to u_\epsilon$ locally
uniformly on $G$.  Thus the weak limit equals $u_\epsilon$ there.
Lemma~\ref{lem:null-area} gives $|B|=0$, so the equality holds almost
everywhere on the sphere.  This proves global Sobolev membership of
$u_\epsilon$, with a bound independent of $\epsilon$.  Finally let
$\epsilon\downarrow0$ and use uniform convergence to $u$ and weak
Sobolev compactness once more.  This proves the theorem.
\end{proof}
